
%% ============================================================
%% ВСТАВКА (23 окт 2026): структурный вывод о планковском секторе G.
%% Вставить после раздела sec:burgers-scale (или перед Заключением).
%% ============================================================

\section{Структурный барьер: гравитации необходим отдельный планковский сектор}
\label{sec:planck-sector}

\subsection{Четыре независимые попытки вывести \texorpdfstring{$b=\ell_P$}{b=lP}}

Тождество $(\ell_P/a)^2=\alpha_G$ (раздел~\ref{sec:burgers-scale})
переводит вопрос о малости гравитации в вопрос <<почему вектор Бюргерса
равен планковской длине, $b=\ell_P$>>. Были проверены четыре независимых
пути вывести это равенство из первых принципов модели, \emph{не} вводя
$G$ заранее:

\begin{enumerate}
\item \textbf{Тождество масштаба.} $(\ell_P/a_{comp})^2=Gm_e^2/\hbar c$ ---
      переписывание определения $\alpha_G$ в планковских единицах; $G$
      уже стоит в $\ell_P$. Круговая зависимость.
\item \textbf{Энергетическое замыкание.} $m_ec^2=\gamma_{3D}b^2a$ с
      $\gamma_{3D}=c^7/8\pi\hbar G^2$ даёт $b\sim10^{-57}$~м
      ($b/\ell_P\to0$); комбинация $\gamma_{3D}\ell_P^2a/m_ec^2$
      воспроизводит $1/\alpha_G$ (расхождение $\sim10^{44}$).
\item \textbf{Условие устойчивости.} Баланс кулоновской самоэнергии и
      линейного натяжения $T=\mu b^2$ с $\mu=\gamma_{3D}$ даёт
      $(b/\ell_P)^2=4\alpha\,\alpha_G$ --- снова с $\alpha_G$ внутри.
\item \textbf{Квантование потока (Брагинский).} Голономия
      $\oint p\,dl=2\pi n\hbar$ квантует $b$ в единицах, построенных из
      $(\hbar,c,e)$, то есть на комптоновском/классическом масштабе
      ($10^{-13}$--$10^{-15}$~м), а не планковском ($10^{-35}$~м).
\end{enumerate}

\subsection{Вывод}

Все четыре пути согласованно не дают $b=\ell_P$: planck-масштаб не
возникает ни из жёсткости вакуума, ни из квантования, построенного на
$(\hbar,c,e)$. Это не случайность отдельного приёма, а структурный факт:
\textbf{квант Бюргерса, фиксируемый электромагнитными и квантовыми
постоянными $(\hbar,c,e)$, несовместим с планковской длиной}. Иными
словами, канал, порождающий кулон (модуль $a_1$ линейной теории), и
канал, который должен был бы порождать ньютоновскую гравитацию с
правильной величиной (модуль $f_1$ на планковском масштабе), управляются
\textbf{физически разными квантами}. Их нельзя связать одним и тем же
вектором Бюргерса.

Отсюда честный статус теории: одновременное описание кулона и гравитации
остаётся \emph{непротиворечивой двухмодульной калибровкой} (два
независимых $a_1,f_1$), но гравитационная постоянная $G$ не выводится ---
для неё требуется \textbf{отдельный, независимо постулируемый планковский
сектор}, вводящий масштаб $\ell_P$ не через $(\hbar,c,e)$. Выявление
этого барьера --- содержательный отрицательный результат: он точно
локализует единственное место, где модель перестаёт быть предсказательной.
